\documentclass[11pt,a4paper]{article}
\usepackage[utf8]{inputenc}
\usepackage[T1]{fontenc}
\usepackage[french]{babel}
\usepackage{amsmath,amssymb,amsthm}
\usepackage{amsfonts}
\usepackage{mathtools}
\usepackage{geometry}
\usepackage{hyperref}
\usepackage{enumitem}
\usepackage{booktabs}
\usepackage{array}
\usepackage{mathrsfs}
\usepackage{xcolor}
\usepackage{microtype}

\geometry{a4paper, top=2.5cm, bottom=2.5cm, left=2.5cm, right=2.5cm}

\theoremstyle{plain}
\newtheorem{theorem}{Th\'eor\`eme}
\newtheorem{lemma}[theorem]{Lemme}
\newtheorem{corollary}[theorem]{Corollaire}

\DeclareMathOperator{\Gal}{Gal}
\DeclareMathOperator{\Hom}{Hom}
% \deg, \det, \dim are already defined in LaTeX
\DeclareMathOperator{\Sw}{Sw}
\DeclareMathOperator{\Ind}{Ind}

\begin{document}

\begin{flushright}
Bures sur Yvette, le 1\textsuperscript{er} juillet 1968
\end{flushright}

\vspace{0.5cm}

Cher Monsieur Weil,

\vspace{0.3cm}

Voici la d\'etermination, dans le cas d'un corps de fonctions, de la fa\c con dont varie la constante de l'\'equation fonctionnelle dont nous avions parl\'e. Je joins une copie d'une lettre de Langlands, et de quelques calculs qui s'y rapportent. Je n'en utiliserai que l'annotation (L) ci-plus bas.

\vspace{0.5cm}

\noindent\textbf{1. Repr\'esentations $\lambda$-adiques, d'apr\`es Serre et Grothendieck.}

\vspace{0.3cm}

Soient $K$ un corps de fonctions d'une variable de corps des constantes $\mathbb{F}_q$, $X$ l'ensemble de ses places, $E$ un corps de nombres, $\lambda$ une place de $E$ ne divisant pas $q$ et $E_\lambda$ le corps local correspondant, d'anneau des entiers $\mathcal{O}_\lambda$ et de corps r\'esiduel $k_\lambda = \mathcal{O}_\lambda/\lambda$.

Une \emph{repr\'esentation $\lambda$-adique} de $K$ est une repr\'esentation continue de $\Gal(\overline{K}/K)$ dans un vectoriel de dimension finie sur $E_\lambda$, strictement c'est une telle repr\'esentation telle qu'il existe une partie finie $S$ de $X$ en dehors de laquelle $\sigma$ soit non ramifi\'e.

Si $v \in X$, on d\'esigne par $V_v$ une uniformisante de $K_v$ correspondante, par $F_v = \varphi_v^{-1}$ un Frobenius g\'eom\'etrique, par $I_v$ le groupe d'inertie du groupe de Galois local en $v$, par $\sigma_v$ la repr\'esentation locale d\'efinie par $\sigma$.

\vspace{0.5cm}

On pose
\[
Z_v(\sigma, T) = \det(1 - F_v T, \sigma_v^{I_v})^{-1}
\]
(action de Frobenius sur les invariants de l'inertie)
\[
Z(\sigma, T) = \prod_v Z_v(\sigma, T) \quad \in E_\lambda[[T]]
\]

\noindent\textbf{Exemple.} Une repr\'esentation $\lambda$-adique de rang $1$ de $K$ n'est autre qu'un homomorphisme du groupe des classes d'id\'eaux de $K$ dans $\mathcal{O}_\lambda^\times$, continu pour la topologie discr\`ete de $\mathcal{O}_\lambda^\times$. On d\'esigne par $1(n)$ la repr\'esentation correspondant \`a $x \mapsto |x|^{-n}$ $(n \in \mathbb{Z})$ et on pose $\sigma(n) = \sigma \otimes 1(n)$.

Un th\'eor\`eme de Grothendieck affirme que si $\sigma$ est une repr\'esentation $\lambda$-adique de $K$, alors

\begin{enumerate}[label=\alph*)]
\item $Z(\sigma, T)$ est une fraction rationnelle (dans $E_\lambda(T)$)
\item Elle v\'erifie une \'equation fonctionnelle
\[
Z(\check{\sigma}(1), T) = \varepsilon(\sigma, T) \, Z(\sigma, T)
\]
o\`u $\varepsilon(\sigma, T)$ est un mon\^ome. Je poserai
\[
\varepsilon(\sigma, T) = q^{(1-g)\dim \sigma} \, \varepsilon_0(\sigma, T)
\]
$\varepsilon_0(\sigma) = \varepsilon_0(\sigma, 1)$ ; cette quantit\'e est une unit\'e $\lambda$-adique.
\item Si $\sigma$ d\'esigne encore le faisceau $\lambda$-adique d\'efini par $\sigma$ (image directe de sa restriction au point g\'en\'erique), on a encore
\[
\varepsilon(\sigma, T) = \det(-FT, H^*(\sigma)) = \prod_i \det(-FT, H^i(\sigma))^{(-1)^i}
\]
\end{enumerate}

\newpage

Si $S$ est une partie finie de $X$, on pose
\[
Z_S(\sigma, T) = \prod_{v \notin S} Z_v(\sigma, T)
\]
et
\[
\varepsilon_S(\sigma, T) = \varepsilon(\sigma, T) \cdot \prod_{v \in S} \det(-F_v T, \sigma_v^{I_v})^{-1}
\]
\begin{align*}
&= \text{mon\^ome asymptotique \`a } Z_S(\sigma, T) \text{ pour } T \to \infty \\
&= \det(-FT, H^*(\sigma_S)) \quad \text{o\`u } \sigma_S \text{ est le faisceau}
\end{align*}
$\lambda$-adique $\sigma$, restreint \`a $X - S$ et prolong\'e par $0$ en les $v \in S$.
\[
\varepsilon_{0,S}(\sigma, T) = q^{(g-1)\dim \sigma} \, \varepsilon_S(\sigma, T)
\]

On d\'efinit les conducteurs
\[
m_v(\sigma_S) = \begin{cases}
\dim(\sigma) - \dim(\sigma_v^{I_v}) + \dim \Hom_{I_v}(\Sw_v, \sigma_v) & \text{si } v \notin S \\[6pt]
\dim(\sigma) + \dim \Hom_{I_v}(\Sw_v, \sigma_v) & \text{si } v \in S
\end{cases}
\]
o\`u on pose par exemple
\[
\dim \Hom_{I_v}(\Sw_v, \sigma_v) = \dim \Hom_{I_v}(\Sw_v, \sigma_v \text{ r\'eduit mod } \lambda)
\]
o\`u $\Sw_v$ est la repr\'esentation de Swan d'un quotient assez grand $I_v'$ de $I_v$ (c'est un $\mathcal{O}_\lambda[I_v']$-module projectif).

On omet la mention de $S$ lorsque $S = \emptyset$.

\vspace{0.5cm}

Soit maintenant $\mathbb{L}$ un ensemble non vide de \emph{premiers} $\neq q$ de places de $E$, et $S$ une partie finie de $X$. Un \emph{syst\`eme compatible} de repr\'esentations $\lambda$-adiques ($\lambda \in \mathbb{L}$), d'ensemble exceptionnel $S$, est la donn\'ee, pour chaque $\lambda \in \mathbb{L}$, d'une repr\'esentation $\lambda$-adique $\sigma_\lambda$ de $K$, les conditions suivantes \'etant v\'erifi\'ees :

\begin{enumerate}[label=(\roman*)]
\item Chaque $\sigma_\lambda$ est non ramifi\'e en dehors de $S$
\item Pour $v \notin S$, $Z_v(\sigma_\lambda, T)$ se trouve dans $E(T)$ et est ind\'ependant de $\lambda \in \mathbb{L}$
\end{enumerate}

On pose
\begin{align*}
Z_v(\sigma, T) &= Z_v(\sigma_\lambda, T) \qquad (v \notin S) \\
Z_S(\sigma, T) &= \prod_{v \notin S} Z_v(\sigma, T) \\
\varepsilon_S(\sigma, T) &= \varepsilon_S(\sigma_\lambda, T)
\end{align*}
ces quantit\'es ne d\'ependent pas de $\lambda$ (en effet, $\varepsilon_S$ peut se d\'efinir en terme de $Z_S$).

\vspace{0.5cm}

\noindent\textbf{2. (L)} \quad On admettra l'\'enonc\'e suivant :

\smallskip

\textit{Soient $\sigma$ et $\rho$ deux repr\'esentations complexes d'un quotient fini de $\Gal(\overline{K}/K)$. Supposons leurs conducteurs disjoints. Avec les notations pr\'ec\'edentes, on a alors}
\begin{equation}\tag{L}
\varepsilon_0(\sigma \otimes \rho, T) = \varepsilon_0(\sigma, T)^{\dim \rho} \cdot \varepsilon_0(\rho, T)^{\dim \sigma} \cdot \prod_{v \in X} \det \sigma(F_v)^{-m_v(\rho)} \cdot \prod_{v \in X} \det \rho(F_v)^{-m_v(\sigma)}
\end{equation}

\newpage

\noindent\textbf{3. Variation des constantes}

\vspace{0.3cm}

\begin{theorem}
Supposons $\mathbb{L}$ infini, et soient $\sigma$ et $\rho$ deux syst\`emes compatibles de repr\'esentations $\lambda$-adiques ($\lambda \in \mathbb{L}$), d'ensembles exceptionnels disjoints. Alors, pour tout $\lambda \in \mathbb{L}$,
\begin{equation}\tag{*}
\varepsilon_0(\sigma_\lambda \otimes \rho_\lambda, T) = \varepsilon_0(\sigma_\lambda, T)^{\dim \rho} \cdot \varepsilon_0(\rho_\lambda, T)^{\dim \sigma} \cdot \prod_{v \in X} \det \sigma_\lambda(F_v)^{-m_v(\rho_\lambda)} \cdot \prod_{v \in X} \det \rho_\lambda(F_v)^{-m_v(\sigma_\lambda)}
\end{equation}
i.e. la formule (L) reste valable.
\end{theorem}

\smallskip

Soient $S$ et $T$ des ensembles exceptionnels disjoints pour $\sigma$ et $\rho$.

\begin{lemma}
Pour $\lambda$ donn\'e, la formule $(*)$ \'equivaut \`a
\begin{equation}\tag{*$_{S,T}$}
\varepsilon_{0,S \cup T}(\sigma_\lambda \otimes \rho_\lambda, T) = \varepsilon_{0,S}(\sigma_\lambda, T)^{\dim \rho} \cdot \varepsilon_{0,T}(\rho_\lambda, T)^{\dim \sigma} \cdot \prod_{v \in T} \det \sigma_\lambda(F_v)^{-m_v(\rho_\lambda)} \cdot \prod_{v \in S} \det \rho_\lambda(F_v)^{-m_v(\sigma_\lambda)}
\end{equation}
\end{lemma}

Si (A) (resp. (B)) est le rapport des deux membres de $(*)$ (resp. de $(*)_{S,T}$), le quotient $(A)/(B)$ appara\^it comme un produit de termes locaux (pour $v \in S \cup T$), dont on va montrer qu'ils sont tous $1$. Si $v \in S$, il faut v\'erifier :
\[
\det(-F_v T, \sigma_v^{I_v} \otimes \rho_v) = \det(-F_v T, \sigma_v^{I_v})^{\dim \rho} \cdot \det \rho(F_v)^{\dim \sigma_v^{I_v}}
\]
ce qui est \'evident.

Si $\sigma$ est une repr\'esentation de $\Gal(\overline{K}/K)$ dans un vectoriel de dimension finie sur $k_\lambda$, on d\'efinit $Z(\sigma, T)$ et $Z_S(\sigma, T)$ dans $k_\lambda[[T]]$ par les m\^emes formules que pr\'ec\'edemment. Il est trivial par r\'eduction mod $\lambda$ que, si $\sigma$ est une repr\'esentation de $\Gal(\overline{K}/K)$ dans un $\mathcal{O}_\lambda$-module libre de type fini, de conducteur contenu dans $S$, et si $\sigma/\lambda$ est la $k_\lambda$-repr\'esentation qui s'en d\'eduit par r\'eduction mod $\lambda$, on a :

\begin{lemma}
\begin{enumerate}[label=(\roman*)]
\item $Z_S(\sigma/\lambda, T)$ est une fraction rationnelle ; on d\'esigne par $q^{(1-g)\dim \sigma} \cdot \varepsilon_0(\sigma/\lambda, T)$ le mon\^ome asymptotique \`a $Z(\sigma/\lambda, T)$ pour $T \to \infty$.
\item Si $\sigma$ et $\rho$ sont comme plus haut, \`a conducteurs contenus dans des parties finies disjointes $S$ et $T$ de $X$, les conditions suivantes sont \'equivalentes :
\begin{enumerate}[label=(\alph*)]
\item $\varepsilon_{0,S \cup T}(\sigma \otimes \rho, T) = \varepsilon_{0,S}(\sigma, T)^{\dim \rho} \cdot \varepsilon_{0,T}(\rho, T)^{\dim \sigma} \cdot \prod_{v \in T} \det \sigma(F_v)^{-m_v(\rho_S)} \cdot \prod_{v \in S} \det \rho(F_v)^{-m_v(\sigma_S)} \pmod{\lambda}$
\item $\varepsilon_{0,S \cup T}(\sigma/\lambda \otimes \rho/\lambda, T) = \varepsilon_{0,S}(\sigma/\lambda, T)^{\dim \rho} \cdot \varepsilon_{0,T}(\rho/\lambda, T)^{\dim \sigma} \cdot \prod_{v \in T} \det \sigma_\lambda(F_v)^{-m_v(\rho_S)} \cdot \prod_{v \in S} \det \rho_\lambda(F_v)^{-m_v(\sigma/\lambda_S)}$ \quad (dans $k_\lambda$)
\end{enumerate}
\end{enumerate}
\end{lemma}

On va d\'eduire de (L) que

\begin{corollary}
\begin{enumerate}[label=(\roman*)]
\item Si $\sigma$ est une repr\'esentation de $\Gal(\overline{K}/K)$ dans un $k_\lambda$-vectoriel de dimension finie, les fonctions $Z_S(\sigma, T) \in k_\lambda[[T]]$ sont rationnelles.
\item Si $\sigma$ et $\rho$ sont deux telles repr\'esentations, \`a conducteurs contenus dans des ensembles finis disjoints $S$ et $T$, la formule (b) est toujours valable.
\end{enumerate}
\end{corollary}

Pour prouver (i), on note que, quitte \`a agrandir $E$ et $k_\lambda$, on sait (Brauer) que $\sigma$ se rel\`eve virtuellement en une repr\'esentation d'un quotient fini de $\Gal(\overline{K}/K)$ dans un vectoriel de dimension finie sur $E$ ; on applique alors 2.(i) \`a cette repr\'esentation.

L'assertion (ii) se prouve de m\^eme, \`a partir de (L) et du lemme 1.

Il en r\'esulte que l'assertion 2.(ii) a) est toujours vraie, donc que $(*)_{S,T}$ est vraie mod $\lambda$.

\newpage

La m\^eme technique de r\'eduction au cas des s\'eries L usuelles, ou une autre m\'ethode, permet de v\'erifier que
\[
\deg_T(\varepsilon_0(\sigma_\lambda, T)) = \dim(\sigma)(2 - 2g) - \sum_v m_v(\sigma_\lambda, S)
\]
(formule d'Ogg--Shafarevitch--Grothendieck). Le premier membre de cette formule est ind\'ependant de $\lambda \in \mathbb{L}$, de sorte que $\sum_v m_v(\sigma_\lambda, S)$ est ind\'ependant de $\lambda \in \mathbb{L}$. Il existe donc une partie infinie $\mathbb{L}_1$ de $\mathbb{L}$ telle que les $m_v(\sigma_\lambda, S)$ (et $m_v(\rho_\lambda, S)$) soient ind\'ependants de $\lambda \in \mathbb{L}_1$. Les deux membres de la formule $(*)_{S,T}$ sont alors ind\'ependants de $\lambda \in \mathbb{L}_1$ ; puisqu'ils sont \'egaux mod $\lambda$ pour une infinit\'e de $\lambda$, ils sont \'egaux pour $\lambda \in \mathbb{L}_1$. Pour achever la d\'emonstration, il suffit de v\'erifier :

\begin{lemma}
Supposons $\mathbb{L}$ infini, et soit $\sigma$ un syst\`eme compatible de repr\'esentations $\lambda$-adiques, d'ensemble exceptionnel $S$. Alors, les entiers $m_v(\sigma_\lambda, S)$ sont ind\'ependants de $\lambda \in \mathbb{L}$.
\end{lemma}

On appliquera les r\'esultats qui pr\'ec\`edent \`a $\sigma$ et \`a un caract\`ere $\chi$ d'ordre fini du groupe des classes d'id\'eaux de $K$, \`a valeurs dans une extension finie de $E$, de conducteur $T$ disjoint de $S$. On sait d\'ej\`a qu'il existe des entiers $m_v$ (les $m_v(\sigma_\lambda, S)$ pour $\lambda \in \mathbb{L}_1$) tels que

$\varepsilon_{0,S \cup T}(\sigma \otimes \chi, T) = \varepsilon_{0,S}(\sigma, T) \cdot \varepsilon_{0,T}(\chi, T)^{\dim \sigma} \cdot \prod_{v \in T} \det \sigma(F_v)^{-m_v(\chi_T)} \cdot \prod_{v \in X} \chi(F_v)^{-m_v}$

et que, d'autre part

$\varepsilon_{0,S \cup T}(\sigma \otimes \chi, T) \equiv \varepsilon_{0,S}(\sigma, T) \cdot \varepsilon_{0,T}(\chi, T)^{\dim \sigma} \cdot \prod_{v \in T} \det \sigma(F_v)^{-m_v(\chi_T)} \cdot \prod_{v \in X} \chi(F_v)^{-m_v(\sigma_\lambda, S)} \pmod{\lambda}$

de sorte que
\[
\prod_{v \in S} \chi(F_v)^{m_v - m_v(\sigma_\lambda, S)} \equiv 1 \pmod{\lambda}
\]
Ceci, \'etant vrai pour tout $\chi$, implique que $m_v(\sigma_\lambda, S) = m_v$, ce qui ach\`eve la d\'emonstration.

\vspace{1cm}

\begin{flushright}
Bien \`a vous

\vspace{0.5cm}

P. Deligne
\end{flushright}

\end{document}
